\begin{abstract}
% Draft; numbers to come from the p2-* grids.
Virtual waiting rooms are sold as the fix for websites that collapse when a
crowd arrives at an announced time, but their behaviour has rarely been
measured. In earlier work we found that a waiting room in front of an
overloaded examination-result portal served almost every student, but that
the obvious designs each fail in a different way: a ticket kept only in the
page sends students who press refresh to the back of the queue, and a ticket
kept in a cookie lets admitted students who run out of patience overload the
portal again, so fewer are served. We design and evaluate three changes: an
admission lease that limits each admitted student to a few requests,
admission whose rate follows the portal's response times instead of a
capacity guessed in advance, and wait estimates shown in the queue. We also
measure how students behave on a slow result page in a study with
\todo{N} students, and use the measurements in the student model of our
testbed. \todo{Headline results: lease, adaptive, estimates, bots.}
\end{abstract}
